\section{Where we start}
\label{sec:start}

Two finished pieces already fit together: Assignment 2's prediction module returns the JSON shapes
that Assignment 1's prototype screens read. Assignment 3 is the layer between them and the features
around them.

\subsection{From Assignment 2: the models and \code{predict.py}}
\label{sec:inherited}

\code{Predictor.load()} reads the saved artefacts once; every method returns plain dictionaries in
camelCase. The timings are from the Mac (Apple M5, CPU environment \code{cos30049-a2}) on
10 November 2026.

{\small
\begin{longtable}{@{}L{4.9cm}L{5.8cm}L{3.3cm}@{}}
\toprule
\textbf{Function} & \textbf{Returns} & \textbf{Measured} \\
\midrule
\endhead
Import and \code{Predictor.load()} & Features, transactions and scores of the snapshot; nine saved methods; the selected method & 4.7 s to import, 0.9 s to load \\
\code{score_address(address)} & Snapshot, range, per-day counts, totals, history status and every transaction's full detail & 0.04 s (273 transactions) to 0.10 s (700); 0.8 to 2.0 MB of JSON \\
\code{transaction(address, txid)} & One transaction's detail & as above (it scores the address) \\
\code{score_what_if(address, amount_btc, direction, time_utc)} & A hypothetical transaction scored against the real, strictly earlier history by eight methods & 0.40 s \\
The same, for many amounts in one feature build & Identical scores for all eight methods & 7 amounts: 0.16 s, against 1.88 s for seven single calls. 41 amounts: 0.18 s \\
\code{explain(tx, use_assistant)} & The template text, or the first Assistant answer that passes the guardrail & under 1 ms for the template; about 10 s for the Assistant on the PC's CPU (README), after a first load of up to a minute \\
\code{txgpt_surprisal(address)} & Surprisal per transaction and per field, from a saved table & a table read \\
\code{classify_heist}, \code{classify_elliptic} & The labelled-data classifiers & not used (W-04) \\
\bottomrule
\end{longtable}
}

\textbf{The data behind them.} The snapshot of 22 October 2026 holds 28 addresses (25 main and 3
with too little history to score) and 7,966 transactions; the largest address has 700 and the
smallest 3. The Explorer's flags come from the robust-z baseline, which flags a transaction when any
feature's $|z|$ reaches 5.0. Nine methods have a saved model and a threshold
(Section~\ref{sec:methods}). \code{Unsupported} and \code{InsufficientHistory} are the exceptions
behind Assignment 1's 404 and 422.

\textbf{What the running example looks like.} Address
\code{1Bbcsfxj2JRUD5crcdZVHBcViT61vrJ3dz} has 377 transactions on 214 days from 16 May 2024 to
19 October 2026; 367 are scored and 3 are flagged, on 3 days. A what-if of 0.05 BTC out at
\code{2026-10-20T12:00:00Z} is flagged with score 5.98: 0.05 BTC against a usual 0.000525 BTC. The
plan uses this address in every example, as Assignment 2's README does.

\subsection{From Assignment 1: the prototype}

\begin{itemize}
  \item A React and TypeScript application built with Vite (\fp{asm1_demo/frontend}, commit
    \code{8375cf0}): about 4,800 lines, 66 tests.
  \item Three linked screens: address selection, daily timeline, day and transaction detail.
  \item Six feedback states, each with a next action: invalid, unsupported, loading, empty,
    insufficient history and network failure.
  \item A hash router, so every screen and state has its own link and the back button works.
  \item Address validation in the browser with the real checksums (Base58Check, Bech32, Bech32m).
  \item \fp{src/api/types.ts}, the contract, and \fp{src/api/http.ts}, a client for the planned
    back end that is switched on with \code{VITE_API_MODE=http}. A mock answers the same five
    endpoints in the browser.
  \item Two charts drawn by hand in SVG, keyboard control of the timeline, a table view of each
    chart, light and dark themes, layouts checked at 1280, 820 and 390 pixels by a screenshot script
    that fails on sideways scrolling.
\end{itemize}

\subsection{The gaps Assignment 3 closes}
\label{sec:gaps}

\begin{center}
\small
\begin{tabularx}{\textwidth}{@{}L{5.2cm}Y L{2.0cm}@{}}
\toprule
\textbf{Today} & \textbf{Needed} & \textbf{Section} \\
\midrule
No server: \code{predict.py} is called from a terminal & A FastAPI service with routes, schemas, validation, errors and documentation & \ref{sec:architecture}, \ref{sec:api} \\
\code{score_address} returns up to 2 MB at once & Small responses per view: timeline, a page of rows, one detail & \ref{sec:api} \\
The prototype shows invented data & The same screens on the real snapshot & \ref{sec:frontend} \\
No form sends user input to a model & The what-if lab, validated on both tiers & \ref{sec:viewwhatif}, \ref{sec:errors} \\
Charts are hand-drawn SVG, two types, no zoom & Plotly.js charts, three or more types, zoom, filters, export & \ref{sec:charts} \\
Nothing updates by itself & Live what-if, replay, live-lookup progress & \ref{sec:realtime} \\
No export, no personal state & CSV, JSON and PNG export; watchlist and preferences & \ref{sec:frontend} \\
The Assistant runs from the command line only & An endpoint with status feedback, a cache and a clear fallback & \ref{sec:assistant} \\
No README, report or video for the web application & Section~\ref{sec:quality} and Section~\ref{sec:report} & \ref{sec:quality}, \ref{sec:report} \\
\bottomrule
\end{tabularx}
\end{center}

\subsection{Loose ends from Assignment 2 that touch this build}
\label{sec:looseends}

\begin{itemize}
  \item \textbf{No \code{v1.0} tag.} The remote has \code{m0}, \code{results-v1} and
    \code{results-v1.1}; the head of \code{from-scratch} is \code{aa56a26} (7 November). The
    submitted commit is tagged before the models are copied (task T02).
  \item \textbf{The usability walkthrough} was moved to Assignment 3 in the report. It is M-15.
  \item \textbf{The blind reviewer check} was still open on 7 November. Could item C-04 gives it a
    place in the application; it is not needed for any Assignment 3 mark.
  \item \textbf{The glossary's fusion sentence} names the old pair (baseline and regression); the
    fusion in use is regression with Isolation Forest. The application never shows glossary text:
    view V5 describes each method from \code{/api/methods}, and free chat with the 152M model is out
    of scope (W-05). No retraining is needed.
  \item \textbf{The Assistant's weights} (610 MB, or about 305 MB as an fp16 export\est{}) are in
    neither git nor, on the Mac, on disk. The GitHub Release was blocked on 7 November because
    \code{gh} is signed in as an account with read access only. S-05 works without the weights and
    says so in the interface.
  \item \textbf{A licence for the team's own code} is still undecided; the zip needs one line in
    \fp{LICENSES.md}.
\end{itemize}
